\chapter{Introduction}\label{chap:intro}
% INTRODUZIONE GENERALE A COSA SONO I SIDE-CHANNEL, IL PERCHè DEL PROBLEMA, IL PERCHè C'è UN PROBLEMA SU AES S-BOX E IL PERCHè SONO STATE INTRODOTTE QUESTE VARIANTI
% RICORDA DI METTERE TUTTE LE CATEGORIE DI SIDE CHANNEL
% CITARE NOSTRO ARTICOLO DELLA TASSONOMIA DEGLI ATTACCHI E VULN HARDWARE (http://ceur-ws.org/Vol-2597/paper-16.pdf)
% CITAZIONI TANTISSIME

The widely known \emph{Spectre}~\cite{kocher2019spectre} and \emph{Meltdown}~\cite{lipp2018meltdown} attacks demonstrated that side-channel analysis can be exploited, in conjunction with other techniques, to target consumer-grade Personal Computers, up to large scale mainframes. Recent Proof of Concepts~\cite{leakypage} demonstrated the feasibility of these attacks using Web technologies like JavaScript, proving the pervasiveness this class of attacks can leverage in real world scenarios. The hardware layer of any modern device plays a primary role in Information System security: it represents, by construction, \emph{the last line of defense} against intrusions~\cite{cosimo2018futuro}. Hardware is directly or indirectly the base all the other layers rely on. Therefore, a hardware-targeting attack may render useless all the defences implemented in the upper layers (like the ones related to \emph{system} and \emph{application software})~\cite{prinetto2020hardware}.

Given the nature of these attacks, it stands to reason that having the broadest possible knowledge of side-channel analysis, including possible countermeasures, is of extreme importance for digital designers, application engineers, and hardware security experts. Indeed, one the main authors in the field of side-channel analysis, Kocher, insists on the importance of having a systematic approach to security at each implementation level: the intercommunication between the various development groups (cipher designers, hardware designer and software engineers) is essential as security faults often involve unanticipated interactions between components designed by different people~\cite{kocher1999differential}.

Over the last decade, academic research confirmed several times the effectiveness and the consequent dangerousness of attacks targeting cryptographic implementations on physical devices. Nowadays, side-channel attacks are considered to be part of cryptanalysis techniques in the same way as linear and differential analysis. Moreover, the convenience that characterizes this type of attacks poses significant risks in the use of IoT devices and embedded systems, especially given the pervasiveness that these have in everyday life. For instance, while linear or differential cryptanalysis needs terabytes of plaintext-ciphertext pairs to break a block cipher, differential power analysis can leverage as just as 2000 bytes of such pairs to break the related physical implementation~\cite{carlet2005highly}. In general, for a malicious actor, the choice of which attack to use is often dictated by how practical the attack is: in the case of an embedded device this unquestionably falls on side-channel attacks.

In recent years, several countermeasures against this class of attacks have been proposed. However, the performance impact that these solutions have on real embedded devices has been proved to be significant and difficult to circumvent. The need for lighter countermeasures has shifted the interest from the implementation/physical level to the logical and mathematical one, the same in which encryption algorithms are specified. The main motivation behind this choice considers resistance to side-channel attacks in the same way as resistance to classical cryptanalysis, treating it as a key parameter to be taken into account for the successful development of old and new block cipher algorithms.

Many of the recent proposals are directly targeting S-Box designs. Since these structures
already provide many properties against classical cryptanalysis, it is therefore reasonable
to integrate in their design possible countermeasures against side-channel attacks. Various publications observed that metrics representing the capability of an \sbocs{} structure to prevent classical cryptanalysis attacks are in opposition with theoretical metrics used to denote the resistance against side-channel analysis.

These observations carried to a slow abandonment of cryptographic design methodologies like algebraic ones, which are able to produce \sbocs{} structures with excellent properties against cryptanalysis but unable to achieve trade offs towards side-channel resistance properties. Therefore, new methodologies with the goal of building novel \sbocs{} structures arose, like the ones based on chaotic maps and heuristic methods. From this point on a myriad of new \sbocs{} proposals were made, claiming of having found a good trade-off among these metrics.

However, we observed that many of these proposal never left the theoretical environment and computer simulations, de facto rarely proving their theoretical claims with everyday embedded implementations. The work carried out in this thesis therefore consists in empirically verifying the resilience to side-channel attacks of some of the latest proposals of \sbocs{} structures. The objective is to test the structures examined in the context of a real attack scenario carried out by means of specialized but low cost and affordable equipment, like the \cw{} board.

\bigskip

\noindent This thesis is organized in five different chapters. 

Chapter~\ref{chap:aes} provides an overview on the \aes{} algorithm and classical cryptanalysis methodologies. Chapter~\ref{chap:pow} gives the necessary basis on side-channel analysis, targeting three of the most powerful and known power attack methodologies. Chapter~\ref{chap:soa} provides an overview on the latest cryptographic proposals targeting \sbocs{} structures, defining the criteria and methodologies used to harden their designs against side-channel attacks. Chapter~\ref{chap:expres} and~\ref{chap:conclusions} finally present the analysis conducted for this thesis work, providing interesting observations on the results obtained and suggesting possible future improvements.